% Source: https://artofproblemsolving.com/community/c3h2008728_2019_state_target_8
% Stable identity: aops-topic-2008728-q01
\textbf{The first phase.} Computing,
\[a_1,a_2,\dots=20,\ 19,\ -1,\ -18,\ 17,\ -1,\ -16,\ 15,\ -1,\ -14,\ 13,\ \dots\]
We claim that for $1\le j\le9$,
\[a_{3j}=-1,\qquad a_{3j+1}=-(20-2j),\qquad a_{3j+2}=19-2j.\]
For $j=1$ this is $a_3=-1$, $a_4=-18$, $a_5=17$. If it holds for $j$ (with $19-2j\ge1$), then
\[a_{3j+3}=|a_{3j+2}|-|a_{3j+1}|=(19-2j)-(20-2j)=-1,\]
\[a_{3j+4}=1-(19-2j)=-(20-2(j+1)),\qquad a_{3j+5}=(18-2j)-1=19-2(j+1),\]
which is the claim for $j+1$. For $j=9$ it gives $a_{27}=-1$, $a_{28}=-2$, $a_{29}=1$, and then $a_{30}=1-2=-1$.

\textbf{The pattern breaks.} Now the absolute values are small and the earlier pattern stops:
\[a_{31}=|a_{30}|-|a_{29}|=1-1=0,\quad a_{32}=0-1=-1,\quad a_{33}=1-0=1,\quad a_{34}=1-1=0,\quad a_{35}=-1,\quad a_{36}=1,\ \dots\]
Since $(a_{34},a_{35})=(a_{31},a_{32})$ and each term depends only on the two before it, the sequence is periodic from $a_{31}$ on with period $3$:
\[a_n=0,\ -1,\ 1\quad\text{for } n\equiv1,\ 2,\ 0\pmod 3,\ n\ge31.\]

\textbf{Answer.} $2019=3\cdot673\equiv0\pmod3$, so $a_{2019}=\boxed{1}$.

(The guess $-1$ comes from assuming that every third term stays $-1$ forever; that is only true up to $a_{30}$.)
